El rendimiento mensual promedio de las empresas con mayor rendimiento acumulado durante el período 2019–2023 evidencia la existencia de un comportamiento estacional relativamente similar entre la mayoría de compañías. Los primeros meses del año (enero, febrero y marzo) registran predominantemente rendimientos promedio negativos, aunque destacan excepciones como
PDD
, que alcanzó un rendimiento promedio de
13.88%
en este período. En contraste,
BNTX
presentó el promedio más bajo, con una caída de
-28.27%
, reflejando una elevada sensibilidad a las condiciones del mercado.
Durante abril y mayo se observa una recuperación generalizada de los rendimientos promedio, donde prácticamente todas las empresas registran valores positivos. Posteriormente, junio representa un punto de desaceleración, con una disminución del rendimiento promedio en la mayoría de compañías. En julio y agosto vuelve a apreciarse un comportamiento favorable, siendo estos algunos de los meses con mejores resultados del período, especialmente para
BNTX
, que alcanzó un rendimiento promedio cercano al
82.76%
.
Finalmente, septiembre y octubre muestran nuevamente una tendencia negativa en la mayoría de empresas, mientras que noviembre y diciembre presentan una recuperación con rendimientos promedio positivos. En términos generales,
BNTX
es la empresa que exhibe la mayor variabilidad entre meses, alternando fuertes caídas y elevados crecimientos, lo que confirma el comportamiento altamente volátil observado en análisis anteriores. Por el contrario, empresas como
FICO
presentan cambios mensuales más moderados, reflejando una mayor estabilidad en sus rendimientos a lo largo del período analizado.